\section{KẾT LUẬN VÀ HƯỚNG PHÁT TRIỂN}\label{sec:ket_luan_huong_phat_trien}

\subsection{Kết quả đạt được}

Đề tài đã xây dựng một mô hình mạng Campus đa cơ sở kết hợp SDN và SD-WAN trên nền tảng EVE-NG, gồm Campus chính (Site 100), ba cơ sở phụ Cần Thơ, Đà Nẵng, Nha Trang (Site 200, 300, 400), trung tâm điều khiển SD-WAN (Site 900), hai nhà cung cấp dịch vụ Internet, MPLS và một cơ sở mở rộng (Site 500). Mô hình có 74 node, 107 phân đoạn mạng, sử dụng thiết bị của nhiều nền tảng: Cisco IOL, vIOS, ASAv, Viptela, Open vSwitch, Ryu, Windows Server, Windows và Ubuntu. Các kết quả chính gồm:

\textbf{Về thiết kế}, đề tài đã xây dựng quy hoạch địa chỉ thống nhất theo mã site, bảng VLAN, bảng liên kết, System-IP, TLOC và ASN cho toàn hệ thống; thiết kế hai kiến trúc cùng tồn tại có chủ đích: ba tầng Core--Distribution--Access có SDN tại Campus chính và Firewall-as-Core tại chi nhánh.

\textbf{Về mạng Campus chính}, hai Core Switch cung cấp gateway VRRP cho năm VLAN, định tuyến OSPF Area~0 với tường lửa biên, DHCP Relay về máy chủ Windows Server; cặp tường lửa ASAv chạy Active/Standby trên các mạng transit dùng chung (VLAN 101--104), chạy OSPF với cả Core và vEdge, tự chuyển vai trò khi lỗi (dưới 1 giây khi chuyển chủ động, khoảng 11--21 giây khi hỏng giao diện hoặc mất hẳn đơn vị Active), bảo vệ vùng DMZ bằng ba danh sách truy cập theo hướng và cho phép quản trị ngoài băng qua ASDM từ VLAN quản trị.

\textbf{Về SDN}, ứng dụng Ryu do nhóm phát triển điều khiển toàn bộ chuyển mạch lớp 2 của sáu Open vSwitch: học MAC theo VLAN, xử lý thẻ 802.1Q theo OpenFlow 1.3, cây chuyển tiếp tính tập trung với thăm dò liên kết chủ động, kênh quản trị VLAN 99 nhiều đường và không phụ thuộc STP, thêm VLAN và áp chính sách bằng một lệnh API, cùng giao diện Campus SDN Console mười trang (tổng quan, topology có cắt/khôi phục liên kết, khôi phục, lưu lượng, VLAN, chính sách có bộ đếm, tải thiết bị, hiệu năng, sức khỏe kênh điều khiển và kịch bản tự động). Kết quả đo cho thấy phần SDN của thao tác thêm VLAN mất 3,8~ms (toàn bộ cấu hình xong trong khoảng 5,3 giây, máy trạm dùng được sau khoảng 40 giây), áp chính sách toàn campus mất 8,0~ms, Controller hội tụ dưới 16~ms, gián đoạn dữ liệu khi mất liên kết không quá 3,5 giây, và cả sáu kịch bản tự phục hồi đều đạt. Dưới tải cao, OVS trong cùng Access đạt khoảng 1,2~Gbit/s, liên VLAN qua Core khoảng 170~Mbit/s, Controller xử lý 1000 MAC mới/giây với 41\% CPU và mặt phẳng điều khiển không bị ảnh hưởng bởi lưu lượng dữ liệu.

\textbf{Về SD-WAN}, mười vEdge được xác thực bằng PKI và hình thành overlay IPsec trên hai transport với BGP underlay; toàn bộ phiên BFD hoạt động; VPN dịch vụ cho phép máy trạm các site thông nhau; vSmart áp tập trung hai SLA class, chính sách AAR cho các lớp lưu lượng, data policy lọc truy cập quản trị từ chi nhánh và control policy ưu tiên vEdge1 tại mọi site; công cụ kiểm thử tự động xác nhận 10/10 phép thử chính sách qua bộ đếm trên vEdge tại cả bốn chi nhánh. Qua 46 lần cắt thật tại bốn chi nhánh, khi mất MPLS, Internet hoặc một vEdge, luồng mới tự khôi phục sau 3,2--13,4 giây nhờ cơ chế theo dõi beacon trên tường lửa chi nhánh; tại Campus chính, sau khi bổ sung phân phối tuyến OMP vào OSPF, mất một transport hoặc mất hẳn vEdge1 được khôi phục sau 7,6--14,4 giây; tắt lần lượt vSmart, vBond, vManage không làm mất gói nào trên overlay; AAR tự chuyển đường khi trễ MPLS vượt ngưỡng SLA. vManage tự gửi thư cảnh báo khi một Edge mất BFD, kết nối điều khiển hoặc OMP.

\textbf{Về truy cập Internet}, người dùng mọi cơ sở ra Internet tại chỗ qua NAT trên vEdge, có dự phòng qua transport MPLS khi mất đường Internet; DNS nội bộ chuyển tiếp tên ngoài; vùng quản trị, Server Farm, DMZ và trung tâm điều khiển chỉ được ra Internet trong phạm vi cần thiết.

\textbf{Về dịch vụ}, hệ thống cung cấp DNS nội bộ \texttt{campus.internal}, Web HTTPS, Mail SMTP/IMAP và Syslog tập trung; người dùng tại Campus chính và tại chi nhánh truy cập được đầy đủ các dịch vụ này; lưu lượng từ chi nhánh đi qua overlay và hai lớp kiểm soát là data policy của SD-WAN và ACL của tường lửa biên.

\textbf{Về vận hành}, nhóm xây dựng Campus NetOps Bot chạy ngay trong topology: người vận hành hỏi đáp và chạy lệnh chỉ đọc qua Telegram, chẩn đoán theo lớp cho vEdge, DHCP, tường lửa và SDN, nhận cảnh báo chủ động từ vManage và SDN Controller (kể cả khi vManage hỏng), và thay đổi cấu hình có duyệt, tự hoàn tác khi trạng thái xấu đi.

\textbf{Về phương pháp}, đề tài hình thành quy trình triển khai từ dưới lên và chẩn đoán theo lớp, có bằng chứng cho từng kết luận và các phép đo được tự động hóa để có thể chạy lại. Quy trình này đã giúp xác định nguyên nhân gốc của nhiều sự cố không hiển nhiên như quy tắc in-band ẩn của Open vSwitch gây bão ARP, trường VLAN khác nhau giữa OpenFlow 1.0 và 1.3, vòng BPDU qua khối OVS làm RSTP rơi vào trạng thái \textit{Dispute}, Controller bị khựng do ONOS chiếm CPU, tuyến cụ thể che mất đường điều khiển của một color TLOC, tuyến mặc định quảng bá ngược làm nhà cung cấp mất Internet, hậu tố \texttt{.local} xung đột với mDNS và giới hạn bảng kết nối của tường lửa không giấy phép.

\subsection{Hạn chế của đề tài}

\begin{itemize}
    \item \textbf{SDN Controller là điểm lỗi đơn:} chỉ có một máy Controller; khi Controller dừng, luồng đã cài vẫn chạy nhưng lưu lượng mới (ARP, DHCP, máy trạm mới) không được xử lý. Khi máy chủ ảo hóa quá tải, Controller khựng tới 4 giây, làm xuất hiện đỉnh trễ và mất gói; chưa xác định được cụ thể tài nguyên nào gây quá tải.
    \item \textbf{Controller SD-WAN chưa có dự phòng:} vManage, vSmart và vBond mỗi loại chỉ có một máy. Mất một controller trong 2 phút không làm mất gói, nhưng trong thời gian đó không thể thay đổi chính sách hay đưa thiết bị mới vào fabric, và mỗi controller cần khoảng 3 phút để kết nối lại sau khi khởi động.
    \item \textbf{\texttt{track} trên tường lửa biên Site 100 chưa được kiểm chứng đầy đủ:} cấu hình \texttt{sla monitor}/\texttt{track} trước đây (theo dõi 10.1.3.2) từng làm cả hai ASAv 9.20 lỗi và khởi động lại liên tục khoảng 1 phút sau khi khởi động, nên đã bị gỡ. Cấu hình mới theo dõi 1.1.1.1 đã qua 23 lần track đổi trạng thái và một lần khởi động lại đơn vị Standby mà không lỗi, nhưng chưa thử khởi động lại đồng thời cả hai đơn vị.
    \item \textbf{Luồng đang chạy khôi phục chậm hơn luồng mới:} khi failover, luồng mới thông lại sau khoảng 3--13 giây nhưng luồng đang chạy mất khoảng 19--38 giây vì phải chờ \texttt{floating-conn} của ASA; khi nối lại liên kết, đôi lúc gián đoạn tới 41 giây. Thời gian phát hiện cũng bị giới hạn bởi timer BFD mặc định, do timer ngắn hơn gây rớt phiên trên vEdge 1 vCPU.
    \item \textbf{Cảnh báo:} thư cảnh báo email phụ thuộc vManage; NetOps Bot bổ sung kênh cảnh báo độc lập nhưng phụ thuộc dịch vụ Telegram và mô hình ngôn ngữ miễn phí có giới hạn số yêu cầu.
    \item \textbf{Giới hạn của môi trường mô phỏng:} thông lượng WAN mô phỏng chỉ khoảng 2~Mbit/s; ASAv không giấy phép giới hạn 100 kết nối và chỉ đạt khoảng 0,1~Mbit/s, nên thông lượng giữa máy trạm các site qua tường lửa không phản ánh thiết bị thật; CPU do Core IOL báo cáo luôn là 0\%, phải đo từ máy chủ; vManage cần tài nguyên lớn; đồng hồ, hiệu năng và hành vi khởi động của thiết bị ảo khác thiết bị thật (ví dụ card mạng ảo đôi khi treo chiều nhận).
    \item \textbf{Một số biện pháp bảo mật mới ở mức thiết kế:} 802.1X, DHCP Snooping, Dynamic ARP Inspection, SNMPv3, SSH thay Telnet toàn hệ thống, xác thực REST API và AAA tập trung.
    \item \textbf{Chính sách SD-WAN cấu hình trực tiếp trên vSmart:} chưa chuyển sang quản lý bằng template và chính sách trên giao diện vManage.
\end{itemize}

\subsection{Hướng phát triển}

\begin{enumerate}
    \item \textbf{Tăng tính sẵn sàng của Controller:} chạy Controller dự phòng nóng (hai tiến trình Ryu với vai trò master/slave của OpenFlow) hoặc chuyển logic sang nền tảng hỗ trợ cluster; tách Controller khỏi các dịch vụ nặng.
    \item \textbf{Hoàn thiện giám sát và cảnh báo:} mở rộng NetOps Bot sang thay đổi cấu hình ASA có hoàn tác và đọc Syslog; tích hợp Syslog và số liệu SNMP vào một bảng điều khiển chung; thêm OSPF graceful restart để rút ngắn thời gian chuyển của tường lửa HA.
    \item \textbf{Giảm thời gian failover:} thử timer BFD ngắn sau khi nâng đủ vCPU cho mọi vEdge, giảm chu kỳ SLA trên Brand-FW và tìm cách để luồng đang chạy không phải chờ \texttt{floating-conn}.
    \item \textbf{Tăng cường bảo mật:} 802.1X tại cổng Access, chống DHCP/ARP giả mạo bằng flow trên Controller, cách ly thiết bị bất thường theo ngưỡng tải, SSH/SNMPv3, xác thực API, AAA tập trung và NTP thống nhất.
    \item \textbf{Chuẩn hóa SD-WAN:} chuyển cấu hình sang template và chính sách quản lý trên vManage.
    \item \textbf{Tự động hóa:} kiểm tra nhất quán giữa tài liệu thiết kế, topology và cấu hình; triển khai và xác minh cấu hình qua API (Ansible, API của vManage, REST của Controller).
    \item \textbf{Đo kiểm trên thiết bị thật:} dựng một cụm nhỏ gồm switch OpenFlow, tường lửa và hai đường WAN thật để đối chiếu hành vi, hiệu năng và thời gian hội tụ với môi trường ảo.
\end{enumerate}

\subsection{Kết luận chung}

Kết quả của đề tài cho thấy SDN và SD-WAN có thể được kết hợp trong một kiến trúc mạng Campus đa cơ sở mà không làm mất vai trò của các cơ chế truyền thống. VRRP, OSPF, tường lửa HA và VLAN tiếp tục bảo đảm độ ổn định tại lớp 3 và biên mạng; Ryu mang lại khả năng lập trình, điều khiển, tự phục hồi và quan sát tập trung cho lớp 2 tại Campus chính, với thời gian thêm VLAN và áp chính sách tính bằng mili giây; còn vManage, vSmart, vBond điều phối overlay an toàn, có chính sách tập trung, tự chọn đường giữa các cơ sở và chủ động cảnh báo sự cố.

Bên cạnh kết quả chức năng, đề tài cũng chỉ ra những thách thức thực tế của SDN in-band --- phụ thuộc vòng giữa switch và Controller, trách nhiệm chống vòng lặp chuyển từ switch sang Controller, tương tác ngầm với STP của phần mạng truyền thống, sự cố dễ lan rộng khi mặt phẳng điều khiển không được thiết kế riêng --- và những yêu cầu chặt chẽ của SD-WAN về chứng thư, định tuyến tới controller, nhãn màu đường truyền và lọc tuyến tại nhà cung cấp. Việc ghi nhận trung thực các giới hạn còn lại là cơ sở để tiếp tục phát triển mô hình thành một nền tảng thực hành, nghiên cứu và thử nghiệm chính sách mạng Campus đa cơ sở.
